\documentclass[]{../../../../tex/classes/styledArticle}
\usepackage{../../../../tex/packages/hebrewSupport}
\usepackage{../../../../tex/packages/mathShortcuts}
\usepackage{../../../../tex/packages/theoremsSupport}

\newcommand\prooftrivial{\begin{proof}[הוכחה באמצעות טרוויאלי]
		טרוויאלי. 
\end{proof}}

\author{שחר פרץ}
\title{\textit{חדו''א 2א $\sim$ הרצאה 1}}
\begin{document}
	\maketitle
	
	\textbf{מרצה: }יעל קרשון
	
	\textbf{חלק מההערות בסיכום זה הן שלי, כלומר הן יכולות להיות שגויות. }
	
	\subsection{נהלים ובירוקרטיה}
	
	חדו''א 2א. יעל קרשון. מייל: yaelkarshon@tauex.tau.ac.il. 
	
	מיילים – 
	\begin{itemize}
		\item דרך המייל האוניברסיטאי
		\item שורת נושא
		\item להתחיל בשלום פרופ' קרסון/שלום יעל
		\item חתימה + שם בסיום
	\end{itemize}
	שעות הקבלה מופיעות במודל. אין צורך לתאם מראש, פשוט לבוא. במשרד שרייבר 126. תנצלו את זה, תבואו לשעות הקבלה של המתרגלים ושאר המרצים. קראו את המשוב על התרגילים. יש חובת הגשה של שמונה תרגילים במהלך הסמסטר. לא הצלחתם? תרגישו תרגיל חלקי. עדיף מכלום. תגישו בזמן, אפילו אם זה רבע תרגיל – באיחור אי אפשר להגיש. 
	
	מילואימניקים; תכתבו, יבואו לקראתכם. בפרט תוכלו להגיש תרגילים באיחור (אם כי הוא לא בהכרח יבדק). 
	
	מקורות עיקריים של הקורס, וכן רשימות של שירי ורשימות של המתרגלים, מופיעות במודל. 
	
	במהלך הסמסטר יהיו שלושה בחנים (''מאוד קצרים``) שיקנו עד 6 נקודות בונוס, בעבור מי שקיבל 56 ומעלה בבחינה הסופית. 
	
	\section{חדו''א 2א}
	אז מה זה חדו''א? חשבון דיפרנציאלי ואינטגרלי. דיפרנציאלי, זה נגזרות, אינטגרלי, זה אינטגרלים. היסטורית, האינטגרלים באו לפני הנגזרות. כבר אצל ארכימדס יש חישובי אינטגרלים מאוד דומים למה שאנו עושים היום (ראו את הדוגמאות של שירי). בחור יווני בשם אפוסטול כתב ספר, גם הוא עם אינטגרלים קודם. החלק הכי חשוב הוא הקשר בין הנגזרות לאינטגרלים. 
	
	נתחיל מהשיעור האחרון ברשימות של שירי – אורך של עקומה בראשון. למה? כי הוא דורש בהתעסקות על המושג של ''חלוקה של קטע``, מושג מאוד חשוב כאשר נדבר על אינטגרלים. באיזשהו הבט, בחדו''א 1א הגדרנו סינוס וקוסינוס גיאומטרית, בלי להסביר מה זה אורך של קשת. זה יסגור את הפינה הזו. 
	
	הפניות לרשימות של סירי יסומנו ב־0.0\{. נעבור על החומרים הבאים: 
	
	2.1\{ אינטגרל מסוים = אינטגרל רימן. \\
	1.2\{ אינטגרל לא מסוים, שהוא קבוצת הפונקציות הקדומות. \\
	1.1\{ חישובים יוריסטיים לאינטגרלים (לקרוא לבד). 
	
	\subsection*{חלוקות של קטע}
	\defi{בהינתן $a, b \in \R$ וכן $a < b$, \textit{קטע סגור} יוגדר להיות $[a, b] := \{c \mid a \le c \le b\}$}
	\exe{בהינתן קטע סגור, ניתן לעבור עליו בקצב אחיד. רמז: 
	\[ \{c \mid a \le c \le b\} = \{(1 - \lg)a + \lg b \mid 0 \le \lg \le 1\} \]}
	\defi{\textit{חֲלוּקָה} של קטע $[a, b]$ (באנגלית – partition) היא קבוצה של $n + 1$ נקודות $\Pi = \{x_0, x_1 \dots x_n\}$ כאשר $n \in \N$, כך ש־$a = x_0 < x_1 < \cdots < x_n = b$. }
	מכאן, שכמות הנקודות בחלוקה היא {סופית}. 
	\defi{בהינתן חלוקה, \textit{תת־הקטע ה־$i$} של החלוקה הוא $[x_{i - 1}, x_i]$. }
	\noti{{אורכו} של תת־הקטע ה־$i$ של החלוקה הוא $\Delta x_i = x_i - x_{i - 1}$. }
	\defi{ממד העדינות של החלוקה הוא $\lg(\Pi) := \max_{i}\Delta x_i$. }
	\defi{\textit{עידון} $\Pi$ של חלוקה $\Pi'$ כך ש־$\Pi' \supset \Pi$. }
	\exe{אם $\Pi' \supset \Pi$ אז $\lg(\Pi') \le \lg(\Pi)$. }\prooftrivial
	עכשיו נעבור לכמה תרגילים יותר מעניינים. 
	\exe{תארו במפורש חלוקה $\Pi$ של $[0, \sqrt 2]$ עם $\lg(\Pi) \le \frac{1}{3}$. }
	\exe{האם $\{0, 1, \frac{1}{2}, \frac{1}{3}, \frac{1}{4}, \dots\}$ היא חלוקה של $[0, 1]$? 
		למה, או למה לא? }
	
	נצטרך לדבר על מושג של מטריקה (שיעול שיעול) מרחק בין נקודות ב־$\R$. לפני כן, נעבור לשעת סיפור. 
	
	חלוקה איננה בהכרח צריכה להיות לקטעים שווים. 
	לפני 25 שנים ממשלת ארה''ב החליטה לעשות מהפכה בלימודי החדו''א בקולג'ים. הם השקיעו מיליונים בלכתוב ספרי מתמטיקה, שיועדו לאנשים שאינם מתמטיקאים וצריכים ללמוד חדו''א. הם הגדירו חלוקה לקטעים שווים, ואינטגרל רימן בהתאם לחלוקה כזו! נראה בעתיד דוגמאות (כמו פונקציית דיריכלה) שבהן ההגדרה הזו נופלת. זו טעות. 
	
	\subsection*{מרחקים ב־$\bm{\R^{n}}$}
	
	\exe{$\sof{a - b} = \sof{a - c} + \sof{c - b}$ אמ''מ $a \le c \le b \lor a \ge c \ge b$}
	''אני מדברת בהמון בטחון עצמי ואומרת טעויות. אתם צריכים לתפוס אותי`` $\sim$ יעל. 
	
	ב־$\R^{n}$ כל נקודה במישור מיוצגת באמצעות קורדינאטות. המרחק בין שתי נקודות $p, q$ מוגדר להיות הנורמה (האוקלידית, לא שזה משנה חדו''אית) של חיסורם. בפרט, ב־$\R^{2}$, בהינתן $p = (x_p, y_p)$ וכן $q = (x_q, y_q)$, נקבל: 
	\[ d(p, q) := \sof{o - q} := \sqrt{(x_p - x_q)^{2} + (y_p - y_q)^{2}} \]
	\exe{
	\[ \sof{x_p - x_q} + \sof{y_p - y_q} \ge \sof{p - q} \ge \begin{cases}
		\sof{x_p - x_q} \\
		\sof{y_p - y_q}
	\end{cases} \]\footnote{רמז: להוכיח את א''ש המשולש בעבור הנורמה האוקלידית באמצעות קושי שוורץ. בעבור המטריקה המושרה מנורמת $L_p$ (מי שיודע מה זה) אפשר להשתמש בא''ש מיניקובסקי במקום, שראינו בחדו''א 1א. }}
	
	\exe{\[ \sof{p - q} = \sof{p - R} + \sof{R - q} \iff R \in [p, q] \]}
	
	\subsection*{עקומות}
	\textbf{לא לבחינה}
		
	\defi{\textit{עקום}/\textit{עקומה}/\textit{מסילה} היא פונקציה $\cg \co [a, b] \to \R^{2}$. }
	זה לא באמת מגביל את הכלליות לדבר על $[0, 1]$ במקום $[a, b]$, בגלל סיבות שנציין בעתיד. טוב לחשוב על $[a, b]$ כעל ''הזמן`` שלאורכו זזה העקומה. 
	
	עכשיו, יש לשאול מהי ההגדרה ההגדרה של רציפות בכלל ב־$\R^{2}$. 
	\begin{Definition}[הגדרה ראשונה לרציפות]
		אפשר לפרק את $\cg$ להיות $\cg(t) = (x(t), y(t))$ (פורמלית: $x = \pi_1 \circ \cg, y = \pi_2 \circ \cg$), ואז לדרוש ש־$x, y$ רציפות. 
	\end{Definition}
	\begin{Definition}[הגדרה שנייה לרציפות]
		לכל $t \in [a, b]$, ולכל $\eg > 0$, קיים $\dg = \dg_{t, \eg} > 0$ כך שלכל $s \in [a, b]$, אם $\sof{t - s} < \dg$ אז $\sof{\cg(s) - \cg(t)} < \eg$. 
	\end{Definition}
	
	יש לשקילות הזו משמעות עמוקה יותר בטופולוגיה\footnote{הדרך העמוקה לראות את זה שההגדרות שקולות, זה כי טופולוגיית טיכונוף היא ה־weak topology ביחס לפונקציות ההיטל $\pi_1 \dots \pi_n$. ואז לא רק שההיטלים רציפים, אלא שזו הטופולוגיה המינימלית (בהינתן טופולוגיה על $\R$) כך שההיטלים רציפים. }. 
	
	\defi{\textit{מסילה לינארית} מ־$p \in \R^{2}$ ל־$q \in \R^{2}$ היא המסילה $\cg = \phi \circ h$, כאשר: 
	\begin{align*}
		[a, b] \rrr{\quad h\quad} &\ [0, 1] \rrr{\quad \phi\quad } \R^{2} \\
		c \mapsto \frac{c - a}{b - a} = &\ \lg \rrr{\quad \ \quad}(1 - \lg)p + \lg q
	\end{align*}
	כאשר $h, \phi$ רציפות. 
	}
	ה־$h$ הזו ש''דוחסת`` את $[a, b]$ ל־$[0, 1]$ היא בדיוק הסיבה שזה לא כזה משנה על איזה domain מסתכלים. \footnote{טופולוגית: הומיהומורפיזם. למעשה יש כאן יותר מזה – הדחיסה מתבצעת ''בקצב קבוע``}
	
	\defi{מסילה פוליגונלית דרך $p_0, p_1 \dots p_n \in \R^{2}$ היא $\cg \co [a, b] \to \R^{2}$ כך שקיימת חלוקה $a = x_0 < x_1 < \cdots < x_n < b$ כך שלכל $i$ מתקיים ש־$\cg|_{[x_{i - 1}, x_i]}$ היא מסילה לינארית מ־$p_{i - 1} = a$ ל־$p_i = b$. }
	
	אינטואיטיבית, זו מסילה שהיא בעצם קווים ישרים (מסילה לינארית) בין נקודות מסוימות במרחב. כמובן שיש להראות שהיא אכן מסילה. 
	
	\noti{למסילה כלשהי $\cg \co [a, b] \to \R^{2}$ ולחלוקה $\Pi = \{x_0 \dots x_n\}$ של הקטע $[a, b]$, נסמן: 
	\[ \ml(\cg, \Pi) = \sum_{i = 1}^{n}\sof{\cg(x_i) - \cg(x_{i - 1})} \]}
	
	\defi{מסילה $\cg \co [a, b] \to \R^{2}$ היא \textit{בעלת אורך סופי} אם הקבוצה: 
	\[ \ccb{\ml(\cg, \Pi) \mid [a, b]\ \textit{חלוקה של}\ \Pi} \]
	חסומה. }
	
	אם כן, 
	\defi{\textit{אורך העקומה} הוא: 
		\[ L(\cg) = \sup\ccb{\ml(\cg, \Pi) \mid [a, b]\ \textit{חלוקה של}\ \Pi} \]
	}
	(במובן הרחב)
	
	\textbf{דוגמה. }נתבונן בעקומה $[-1, 1] \to \R^{2}$ המוגדרת ע''י: 
	\[ \cg(x) = (x, \sqrt{1 - x^{2}}) \]
	\theo{$\cg$ היא עקומה בעלת אורך סופי, וכן $2 \le L(\cg) \le 4$. זו ממש ההגדרה}\begin{proof}
		מספיק להראות:
		\begin{enumerate}
			\item קיימת חלוקה $\Pi$ של $[-1, 1]$ עם $\ml(\cg, \Pi) = 3$
			\item לכל חלוקה $\Pi$ של $[-1, 1]$ עם $\ml(\cg, \Pi) \le 4$
		\end{enumerate}
		פרטים יעלו למודל. 
	\end{proof}
	\defi{$\pi := L(\cg)$ כאשר $\cg$ העקומה המוגדרת לעיל. }
	
	\subsection*{סכום רימן}
	נתחיל מתיאור יוריסטי: ''שטח מתחת לגרף``. עבור פונקציה $f \co [a, b] \to \R$, נסמן: 
	\[ \int^{b}_a \dequad \ f = \int^{b}_a \dequad\ f(x)\dx = \int^{b}_a \dequad \ f(t) \dt = \int f \]
	כאשר ''מתחת לגרף`` השטח שלילי. זכרו לקרוא את סעיף 1.1 ברשימות של שירי לבד, וזה מושג שנקרא ''אינטגרל לא אמיתי`` (זה לא בדיוק אותו הדבר כמו אינטגרל רימן). נתארו גם בצורה יוריסטית; נתבונן בפונקציה $f(x) = \frac{1}{x^2}$ בקטע $f \co [1, \inft) \to \R$. יש לנו כאן בעיה – הפונקציה מוגדרת על קטע (קרן) לא חסום. עתה נתבונן בפונקציה $f(x) = \frac{1}{\sqrt x} = x^{-\frac{1}{2}}$, שמוגדרת על קטע סופי $f \co (0, 1] \to \R$. כאן אומנם הפונקציה מוגדרת על קטע סופי, אך ב־$0$ היא שואפת לאינסוף, גם אם השטח לכאורה סופי. הפתרון הוא לקחת גבול על אינטגרל רימן של הפונקציה. 
	
	\exe{מה הקשר בין השטחים הנ''ל? }
	
	\exe{נתון מלבן $a \times b$, מחולק למלבנים $a_i \times b_i$, כאשר $i \in [n]$. אז: 
	\[ ab = \sum_{i = 1}^{n}a_ib_i \]}
	למה התרגיל הזה חשוב? כי אם נתבסס על הקונספט של ''שטח מלבן`` (שהוא ידוע להיות $a \cdot b$, אורך כפול רוחב) חשוב לראות שזה מכבד את מה שהיינו ''מצפים`` מהגדרה של שטח – שהשלם שווה לסך חלקיו. 
	
	להלן הרעיון מאחורי אינטגרל רימן: בהינתן פונקציה וגרפה, [איור שלא הכנסתי], נבחר חלוקה של הקטע, ''נדגום`` את הערך של הפונקציה בנקודה שבחרנו, וניקח את הערך בנקודת הדגימה להיות הגובה של המלבן. לבחירה הזו נקרא ''תיוג`` או ''נקודות מתאימות``. 
	
	\defi{\textit{חלוקה מתויגת} (באנגלית tagged partition, ברשימות ''חלוקה+נקודות מתאימות``) של הקטע $[a, b]$
		היא זוג $(\Pi, \{t_i\})$ כאשר $\Pi\co a = x_0 < x_1 < \cdots < x_n = b$ חלוקה של הקטע $[a, b]$, וכן תיוג $t_i \in [x_{i - 1}, x_i]$ לכל אחד מ־$i \in [n]$ קטעי החלוקה. אז \textit{סכום רימן} של הפונקציה $f \co [a, b] \to \R$ ביחס לחלוקה המתויגת $\{t_i\}_{i = 1}^{n}$ הוא: 
		\[ S(f, \pi, \{t_i\}) := \sumnio f(t_i)\Delta x_i \]}
		
	\textbf{דוגמה. }(גרף)
	
	\exe{ציירו חלוקה מתויגת של הקטע $[0, 1]$ עם שלושה תתי־קטעים באורכים שאינם כולם שווים, ועם $t_2 = t_3$. }
	
	עתה נתעסק עם שתי דוגמאות \textit{חשובות} שבהחלט יכולות להופיע בבחינה. נתבונן בפונקציה: 
	\begin{align*}
		f \co [a, b] \to \R && f(x) = \begin{cases}
			5 & x = a \\
			15 & x \neq a
		\end{cases}
	\end{align*}
	נתבונן בסכום רימן ביחס לחלוקה כלשהי: 
	\[ S(f, \Pi, \{t_i\}) = \sum_{i = 1}^{n}f(t_i)\Delta x_i = f(t_1)\Dg x_1 + \sum_{i = 2}^{n}\underbrace{f(t_i)}_{5}\Delta x_i \]
	כתלות בבחירה של נקודת הדגימה $t_1$, אם $t_1= a$ אז הסכום יוצא $10 \Delta x_i + 5(b - a)$, אחרת זה יוצא $5(b - a)$. לכן, נקבל: 
	\[ \sof{S(f, \Pi, \{t_i\}) - 5(b - a)} \le 10 \Dg x_i \le 10 \, \lg(\Pi) \]
	
	עתה נסתכל על דוגמה נוספת. 
	
	\textbf{דוגמה. }נתבונן בפונקציית דיריכלה: 
	\[ D(x) = \begin{cases}
		1 & x \in \Q \\
		0 & x \notin \Q
	\end{cases} \]
	ניקח חלוקה מתויגת $(\Pi, \{t_i\})$ של קטע $[a, b]$. 
	\begin{itemize}
		\item אם לכל $i \in [n]$ מתקיים $t_i \in \Q$ (בהכרח קיים תיוג כזה, מצפיפות), אז:
		\[ S(f, \Pi, \{t_i\}) = \sumnio 1 \Delta x_i = b - a \]
		\item אם לכל $i \in [n]$ מתקיים $t_i \notin \Q$ אז: 
		\[ S(f, \Pi, \{t_i\}) = \sumnio 0 \Delta x_i = 0 \]
	\end{itemize}
	
	
	\subsection*{אינטגרל רימן}
	\defi{נתונה פונקציה $f \co [a, b] \to \R$. נאמר ש־$f$ \textit{\href{https://eincyclopedia.org/wiki/\%D7\%A4\%D7\%95\%D7\%A0\%D7\%A7\%D7\%A6\%D7\%99\%D7\%94\_\%D7\%90\%D7\%99\%D7\%A0\%D7\%98\%D7\%92\%D7\%A8\%D7\%91\%D7\%99\%D7\%9C\%D7\%99\%D7\%AA}{אינטגרבילית} רימן עם אינטגרל שווה ל־$I$} אם\footnote{בהגדרות, ''אם`` זה ''אם ורק אם``, כי ברור ששני הכיוונים נדרשים} לכל $\eg > 0$ קיימת $\dg > 0$  כך שלכל חלוקה מתויגת $(\Pi, \{t_i\})$ של $[a, b]$, עם מדד עדינות $\lg(\Pi) < \dg$, מתקיים: 
	\[ \sof{S(f, \Pi, \{t_i\}) - I} < \eg \]
	אם זה קורה, נאמר ש־$I$ \textit{האינטגרל המסויים} של $f$ בקטע $[a, b]$. }
	
	במילים אחרות, לכל אפסילון, עם מספר מסוים לעידון, אז כל חלוקה ותיוג יהיו בסביבת $\eg$ של ערך האינטגרל. 
	
	\noti{מסמנים $\int^{b}_af = \int^{b}_af(x)\dx = I$}
	\defi{הפונקציה $f$ נקראת \textit{אינטגרבילית} אם קיים $I$ כך ש־$\int^{b}_a f = I$. }
	
	\rmark{$I$ איננו במובן הרחב, כלומר $I \in \R$. }
	
	\begin{Exercise}[''לא יכולה  להגזים בכמה הוא חשוב``]
		רשמו את השלילה של הביטוי $\int^{b}_a f = I$. 
	\end{Exercise}
	
	\rmark{נניח שהאינטגרל של $f$ איננו שווה ל־$I$. יש שתי סיבות, מאוד שונות, לסיבה שזה קורה: 
	\begin{itemize}
		\item $\int^{b}_a f = I' \neq I$
		\item $f$ איננה אינטגרבילית רימן בקטע $[0, 1]$. 
	\end{itemize}}
	בהמשך נראה קריטריונים לאינטגרביליות של דורשים לדעת מי זה בדיוק $I$ (באופן דומה לכך שיודעים שפונקציה מתכנסת לאנשהו באמצעות קריטריון קושי, אך לא יודעים לאן). 
	
	נציין את המשפטים הייסודים של החדו''א. לא נוכיח אותם כרגע, הם בחומר לעוד איזה חודש, אבל הם מה שיקשור את התיאוריה של ההרצאה לשבועיים הראשונים של התרגול (שכרגע יראו לא קשורים בעליל). ידועים גם בתור משפטי ניוטון־לייבניץ\footnote{המלחמה בנוגע מי הגיע למשפטים קודם, נמשכת עד היום. }
	
	בהינתן $f, F \co [a, b] \to \R$: 
	\begin{enumerate}
		\item ''תחת הנחות מסוימות`` (\href{https://eincyclopedia.org/wiki/\%D7\%A4\%D7\%95\%D7\%A0\%D7\%A7\%D7\%A6\%D7\%99\%D7\%94\_\%D7\%A1\%D7\%99\%D7\%9E\%D7\%A4\%D7\%98\%D7\%99\%D7\%AA}{פונקציה סימפטית מספיק}), אם $F(x) = \int^{x}_a f(t) \dt$, אז $F'(x) = f(x)$. 
		\item ''תחת הנחות מסוימות``, אם $F'(x) = f(x)$, אז $F(x) = F(a) + \int^{x}_af(t)\dt$
	\end{enumerate}
	זה מסביר למה המושג הבא כל־כך חשוב לחישובי אינטגרלים: 
	\defi{לפונקציות ממשיות $f, F \co I \to \R$ על קטע $I \subset \R$, אז $F$ נקראת \textit{פונקציה קדומה של $f$ על קטע\footnote{לשימושנו, קבוצה סגורה. } $I$} אם $\forall x \in I \co F'(x) = f(x)$, כאשר בקצוות הקטע (אם יש כאלו) מדובר על נגזרת חד־כיוונית. }
	נעבוד בתרגולים למצוא פונקציות קדומות, כי הן בתורן יאפשרו לחשב אינטגרלים מסוימים, באמצעות המשפט היסודי. 
	
	\defi{\textit{האינטגרל הלא מסויים של $f$} הוא \textbf{קבוצת} הפונקציות הקדומות של $f$. }
	בקורס ''פונקציות ממשיות`` נגדיר מושג אינטגרל רחב יותר, שנקרא \textit{אינטגרל לבג}, וניתן להגדיר אותו גם מחוץ ל־$\R$. כרגע אנחנו עובדים עם אינטגרל רימן. 
	
\end{document}